\documentclass[../Odyssey_Project.tex]{subfiles}

\begin{document}
\setcounter{section}{9}
\section{Composition Series}

\subsubsection*{Composition Series}
\begin{itemize}
  \item A series of subgroups $G = G_0 > G_1 > \ldots > G_r = 1$ of a group $G$ is a \emph{composition series} of $G$ if $G_{i+1} \trianglelefteq G_i$ for every $i$ and if each \emph{successive quotient} $G_i / G_{i+1}$ is simple.
  \item The composition series above is said to have \emph{length} $r$.
  \item The successive quotients of a composition series are called the \emph{composition factors} of the series.
  \item More generally, a group is said to be a \emph{composition factor} of $G$ if it is isomorphic with one of the composition factors in some composition series of $G$.
  \item Example: $S_5 > A_5 > 1$ is the only composition series of $S_5$ (since $A_5$ is simple by Theorem~\ref{thm:7.4} and is the only non-trivial proper normal subgroup of $S_5$ by Theorem~\ref{thm:7.5}), with composition factors $\Z_2$ and $A_5$.
\end{itemize}

\subsubsection*{Maximal Normal Subgroups}
\begin{itemize}
  \item $N$ is a \emph{maximal normal subgroup} of $G$ if $N \trianglelefteq G$ and there is no proper normal subgroup of $G$ that properly contains $N$.
  \item By \nameref{thm:correspondence}, $N$ is a maximal normal subgroup of $G$ iff $G/N$ is simple.
  \item A maximal subgroup that is also normal is clearly a maximal normal subgroup (and must be of prime index), but a maximal normal subgroup need not be a maximal subgroup, for it could be properly contained in a proper subgroup that is not normal.
  \item For example, $A_5 \times \Z_2$ has $1 \times \Z_2$ as a maximal normal subgroup that is not maximal.
\end{itemize}

\begin{proposition}
\label{prop:10.1}
Finite groups have composition series.
\end{proposition}
\begin{proof}
We use induction on $\lvert G\rvert$. If $G=1$, then the series $G=1$ is a composition series of length $0$. If $G$ is non-trivial and simple, then $G>1$ is a composition series of $G$. Suppose now that $G$ is not simple. Since $G$ is finite, the set of proper normal subgroups of $G$ has a maximal element; call it $G_1$. Then $G_1$ is a maximal normal subgroup of $G$, so $G/G_1$ is simple by the \nameref{thm:correspondence}. Also $\lvert G_1\rvert<\lvert G\rvert$, so the induction hypothesis gives a composition series
\[
  G_1>G_2>\cdots>G_r=1
\]
of $G_1$. Therefore
\[
  G=G_0>G_1>G_2>\cdots>G_r=1
\]
is a composition series of $G$: the first quotient $G/G_1$ is simple, and the remaining successive quotients are simple by the composition series of $G_1$.
\end{proof}

\begin{itemize}
  \item Infinite groups need not have composition series. For example, using Theorem~\ref{thm:1.5} we see that every non-trivial subgroup of the infinite cyclic group $\Z$ is isomorphic with $\Z$; as $\Z$ is not simple, $\Z$ has no simple subgroups, and hence we cannot construct a composition series of $\Z$, since the last non-trivial term of such a series must be a simple subgroup of $\Z$.
\end{itemize}

\begin{lemma}
\label{lem:10.2}
Let $G$ be a group that has a composition series, and let $N \trianglelefteq G$. Then $N$ has a composition series.
\end{lemma}
\begin{proof}
Let
\[
  G=G_0>G_1>\cdots>G_r=1
\]
be a composition series of $G$, and put $N_i=N\cap G_i$ for each $i$. Then
\[
  N=N_0\geq N_1\geq\cdots\geq N_r=1
\]
is a series of subgroups of $N$. For each $i$, we have $N_{i+1}\trianglelefteq N_i$, because $G_{i+1}\trianglelefteq G_i$ and $N_i\leq G_i$. Moreover,
let $\psi_i\colon G_i\to G_i/G_{i+1}$ be the natural map. Then the restriction of $\psi_i$ to $N_i$ has kernel $N_i\cap G_{i+1}=N_{i+1}$, so by \nameref{thm:first-iso} we have
\[
  N_i/N_{i+1}
  \cong
  \psi_i(N_i)
  =
  (N\cap G_i)G_{i+1}/G_{i+1}.
\]
The subgroup $(N\cap G_i)G_{i+1}/G_{i+1}$ is normal in $G_i/G_{i+1}$, since $N\trianglelefteq G$ gives $N\cap G_i\trianglelefteq G_i$. Since $G_i/G_{i+1}$ is simple, it follows that either $N_i/N_{i+1}$ is trivial or $N_i/N_{i+1}\cong G_i/G_{i+1}$ is simple. Deleting any repeated terms from
\[
  N=N_0\geq N_1\geq\cdots\geq N_r=1
\]
therefore gives a composition series of $N$.
\end{proof}

\subsubsection*{Equivalence of Composition Series}
\begin{itemize}
  \item Let $G = G_0 > G_1 > \ldots > G_r = 1$ be a composition series, and suppose $G = H_0 > H_1 > \ldots > H_r = 1$ is another composition series of the same length $r$. These series are \emph{equivalent} if there is some $\rho \in S_r$ such that $G_{i-1}/G_i \cong H_{\rho(i)-1}/H_{\rho(i)}$ for every $i$.
  \item Example: $G = \langle x \rangle \cong \Z_6$ with $G_1 = \langle x^2 \rangle$, $H_1 = \langle x^3 \rangle$ gives equivalent composition series $G > G_1 > 1$ and $G > H_1 > 1$, since $G/G_1 \cong H_1/1 \cong \Z_2$ and $G_1/1 \cong G/H_1 \cong \Z_3$. (Here $\rho = (1\ 2) \in S_2$.)
  \item Up to equivalence, a group has at most one composition series, so a group having a composition series has a well-defined collection of composition factors.
  \item Since composition factors are simple groups, classifying finite simple groups (completed in 1980) gives comprehensive understanding of finite groups.
\end{itemize}

\begin{namedtheorem}[Jordan--H\"older Theorem]
\label{thm:jordan-holder}
Suppose that $G$ is a group that has a composition series. Then any two composition series of $G$ have the same length and are equivalent.
\end{namedtheorem}
\begin{proof}
Let
\[
  G=G_0>G_1>\cdots>G_r=1
\]
and
\[
  G=H_0>H_1>\cdots>H_s=1
\]
be two composition series of $G$. We use induction on $r$. If $r=0$, then $G=1$, and there is nothing to prove. If $r=1$, then $G$ is simple, so $G>1$ is the only composition series of $G$. Thus assume $r>1$, and assume the result is known for every group having a composition series of length less than $r$.\\
If $G_1=H_1$, then the two tails
\[
  G_1>G_2>\cdots>G_r=1
\]
and
\[
  H_1>H_2>\cdots>H_s=1
\]
are composition series of $G_1$. By induction they have the same length and are equivalent, so the original two composition series of $G$ also have the same length and are equivalent.\\
We may therefore assume that $G_1\neq H_1$. Since $G_1\trianglelefteq G$ and $H_1\trianglelefteq G$, Proposition~\ref{prop:1.7} gives $G_1H_1\trianglelefteq G$. Also $G_1H_1\neq G_1$, for otherwise $H_1\leq G_1$, contradicting the maximal normality of $H_1$ in $G$. Since $G/G_1$ is simple and $G_1<G_1H_1\trianglelefteq G$, it follows that $G_1H_1=G$. Put $K=G_1\cap H_1$. By \nameref{thm:first-iso}, applied to the natural maps into $G/G_1$ and $G/H_1$, we have
\[
  H_1/K\cong G/G_1
  \qquad\text{and}\qquad
  G_1/K\cong G/H_1.
\]
In particular, $H_1/K$ and $G_1/K$ are simple. Since $K\trianglelefteq G$, Lemma~\ref{lem:10.2} gives a composition series
\[
  K=K_0>K_1>\cdots>K_t=1
\]
of $K$, with $t=0$ if $K=1$.\\
Now compare two composition series of $G_1$:
\[
  G_1>G_2>\cdots>G_r=1
\]
and
\[
  G_1>K>K_1>\cdots>K_t=1.
\]
By induction these have the same length and are equivalent. Hence $t+1=r-1$, so $t=r-2$. Similarly, compare the two composition series of $H_1$:
\[
  H_1>H_2>\cdots>H_s=1
\]
and
\[
  H_1>K>K_1>\cdots>K_{r-2}=1.
\]
Induction gives $s-1=r-1$, so $s=r$, and these two series are equivalent. Finally, the two intermediate composition series
\[
  G>G_1>K>K_1>\cdots>K_{r-2}=1
\]
and
\[
  G>H_1>K>K_1>\cdots>K_{r-2}=1
\]
are equivalent, because $G/G_1\cong H_1/K$, $G/H_1\cong G_1/K$, and all remaining factors are the same. Combining these equivalences shows that the original two composition series of $G$ are equivalent and have the same length.
\end{proof}

\subsubsection*{Subnormal, Normal, and Chief Series}
\begin{itemize}
  \item A series of subgroups $G = G_0 \geq G_1 \geq \ldots \geq G_r = 1$ of a group $G$ is called a \emph{subnormal series} of $G$ if $G_{i+1} \trianglelefteq G_i$ for every $i$.
  \item A subnormal series is called a \emph{normal series} if $G_i \trianglelefteq G$ for every $i$.
  \item Some authors use ``normal series'' for what we call a subnormal series, and ``invariant series'' for what we call a normal series.
  \item Composition series are examples of subnormal series, but subnormal series can have successive quotients that are trivial or that are non-trivial but not simple.
  \item Two subnormal series of the same length are said to be \emph{equivalent} if they satisfy the condition stated as the definition of equivalence of composition series.
  \item A \emph{refinement} of a subnormal series is one obtained by interposing additional terms; it is \emph{proper} if at least one of the interposed terms was not already present.
  \item A composition series is a subnormal series that has no repeated terms and does not admit a proper refinement.
  \item A \emph{chief series} of $G$ is a normal series of $G$ with no repeated terms and with the additional property that no normal subgroup of $G$ is contained properly between any two terms of the series.
  \item Analogy: A composition series is a subnormal series having no subnormal series as a proper refinement, and a chief series is a normal series having no normal series as a proper refinement.
  \item The successive quotients of a chief series are its \emph{chief factors}; a group is a \emph{chief factor} of $G$ if it is isomorphic with a chief factor in some chief series of $G$.
  \item The analogue of the \nameref{thm:jordan-holder} for chief series holds---namely, any two chief series of a group have the same length and are equivalent, with a proof virtually identical to that for composition series. This is a special case of the Jordan--H\"older theorem for groups with operators (see [26, Section 2.3]).
\end{itemize}

\subsubsection*{Minimal Normal Subgroups}
\begin{itemize}
  \item $N$ is a \emph{minimal normal subgroup} of $G$ if $1 \neq N \trianglelefteq G$ and there is no non-trivial normal subgroup of $G$ that is properly contained in $N$.
  \item Any non-trivial finite group has minimal normal subgroups.
  \item A simple group has a unique minimal normal subgroup, namely itself.
\end{itemize}

\begin{proposition}
\label{prop:10.3}
Finite groups have chief series.
\end{proposition}
\begin{proof}
We use induction on $\lvert G\rvert$. If $G=1$, then the series $G=1$ is a chief series of length $0$. If $G$ is non-trivial and simple, then $G>1$ is a chief series of $G$. Suppose now that $G$ is not simple. Since $G$ is finite, it has a proper minimal normal subgroup $N$. By induction, the finite group $G/N$ has a chief series. By the \nameref{thm:correspondence}, this chief series has the form
\[
  G/N=G_0/N>G_1/N>\cdots>G_r/N=1,
\]
where $G_i\trianglelefteq G$ for every $i$, $G_0=G$, and $G_r=N$. Moreover, no normal subgroup of $G$ lies properly between $G_{i-1}$ and $G_i$, for such a subgroup would correspond to a normal subgroup of $G/N$ lying properly between $G_{i-1}/N$ and $G_i/N$. Since $N$ is minimal normal in $G$, no non-trivial normal subgroup of $G$ lies properly between $N$ and $1$. Therefore
\[
  G=G_0>G_1>\cdots>G_r=N>1
\]
is a chief series of $G$.
\end{proof}

\begin{itemize}
  \item By definition, the composition factors of finite groups are simple groups. We now turn to determining the nature of chief factors of finite groups; the first step is to rephrase the problem in different terms.
\end{itemize}

\begin{lemma}
\label{lem:10.4}
Let $G$ be a group having a chief series. Then every chief factor of $G$ is a minimal normal subgroup of a quotient group of $G$.
\end{lemma}
\begin{proof}
Let
\[
  G=G_0>G_1>\cdots>G_r=1
\]
be a chief series of $G$. A chief factor of this series has the form $G_i/G_{i+1}$ for some $i$. Since $G_i\trianglelefteq G$ and $G_{i+1}\trianglelefteq G$, the quotient $G_i/G_{i+1}$ is a normal subgroup of $G/G_{i+1}$. If there were a non-trivial normal subgroup of $G/G_{i+1}$ properly contained in $G_i/G_{i+1}$, then by the \nameref{thm:correspondence} it would have the form $M/G_{i+1}$ for some normal subgroup $M\trianglelefteq G$ satisfying
\[
  G_{i+1}<M<G_i.
\]
This contradicts the definition of a chief series. Hence $G_i/G_{i+1}$ is a minimal normal subgroup of the quotient group $G/G_{i+1}$. Therefore every chief factor of $G$ is a minimal normal subgroup of a quotient group of $G$.
\end{proof}

\begin{theorem}
\label{thm:10.5}
A minimal normal subgroup of a finite group is a direct product of mutually isomorphic simple groups.
\end{theorem}
\begin{proof}
Let $G$ be a finite group, and let $N$ be a minimal normal subgroup of $G$. Choose a maximal normal subgroup $N_1$ of $N$, so $N/N_1$ is simple. Let $N_1,N_2,\dots,N_r$ be the distinct conjugates of $N_1$ in $G$. Since $N\trianglelefteq G$, each $N_i$ is a maximal normal subgroup of $N$, and conjugation by an element of $G$ gives an isomorphism $N/N_1\cong N/N_i$. Hence all quotients $N/N_i$ are mutually isomorphic simple groups.\\
Conjugation by elements of $G$ permutes the set $\{N_1,\dots,N_r\}$, so
\[
  I=N_1\cap N_2\cap\cdots\cap N_r
\]
is normal in $G$. Since $I\leq N_1<N$, the minimality of $N$ gives $I=1$.\\
For $1\leq i\leq r$, put
\[
  A_i=N_1\cap\cdots\cap N_i.
\]
We prove by induction on $i$ that $N/A_i$ is a direct product of groups isomorphic with $N/N_1$. For $i=1$, this is immediate. Suppose $i>1$ and the claim holds for $i-1$. If $A_{i-1}\leq N_i$, then $A_i=A_{i-1}$, and there is nothing new to prove. Otherwise $N_i<A_{i-1}N_i\trianglelefteq N$. Since $N_i$ is maximal normal in $N$, we have $A_{i-1}N_i=N$. By Lemma~\ref{lem:2.7},
\[
  N/A_i
  \cong
  A_{i-1}/A_i \times N_i/A_i.
\]
Now the natural map $A_{i-1}\to N/N_i$ has kernel $A_{i-1}\cap N_i=A_i$ and image $A_{i-1}N_i/N_i=N/N_i$, so by \nameref{thm:first-iso},
\[
  A_{i-1}/A_i\cong N/N_i\cong N/N_1.
\]
Similarly, the natural map $N_i\to N/A_{i-1}$ has kernel $N_i\cap A_{i-1}=A_i$ and image $A_{i-1}N_i/A_{i-1}=N/A_{i-1}$, so by \nameref{thm:first-iso},
\[
  N_i/A_i\cong N/A_{i-1}.
\]
By the induction hypothesis, $N/A_{i-1}$ is a direct product of groups isomorphic with $N/N_1$. Therefore $N/A_i$ is also a direct product of groups isomorphic with $N/N_1$. Taking $i=r$ gives $A_r=I=1$, and hence $N$ itself is a direct product of mutually isomorphic simple groups.
\end{proof}

\begin{corollary}
\label{cor:10.6}
A chief factor of a finite group is a direct product of mutually isomorphic simple groups.
\end{corollary}
\begin{proof}
Let $G_i/G_{i+1}$ be a chief factor of a finite group $G$. By Lemma~\ref{lem:10.4}, this chief factor is a minimal normal subgroup of the quotient group $G/G_{i+1}$. Since $G/G_{i+1}$ is finite, Theorem~\ref{thm:10.5} applies to this quotient group. Therefore $G_i/G_{i+1}$ is a direct product of mutually isomorphic simple groups.
\end{proof}

\subsection*{Exercises}

\begin{exercise}
\label{ex:10.1}
Show that an abelian group has a composition series iff it is finite.
\end{exercise}
\begin{solution}
\end{solution}

\begin{exercise}
\label{ex:10.2}
Show that $\mathrm{GL}(n, F)$ has a composition series iff $F$ is finite.
\end{exercise}
\begin{solution}
\end{solution}

\begin{exercise}
\label{ex:10.3}
Using \nameref{thm:third-iso} (Exercise~\ref{ex:2.14}), prove the \emph{Schreier refinement theorem}: Any two subnormal series of a given group have equivalent refinements.
\end{exercise}
\begin{solution}
\end{solution}

\begin{exercise}
\label{ex:10.4}
(cont.) Use the Schreier refinement theorem to give an alternate proof of \nameref{thm:jordan-holder}.
\end{exercise}
\begin{solution}
\end{solution}

\begin{exercise}
\label{ex:10.5}
Determine all composition series and all chief series of $S_n$ for $n \geq 2$. (\textsc{Hint}: Use Exercise~\ref{ex:3.8}.)
\end{exercise}
\begin{solution}
\end{solution}

\begin{exercise}
\label{ex:10.6}
Show that any group having a composition series has a chief series.
\end{exercise}
\begin{solution}
\end{solution}

\begin{exercise}
\label{ex:10.7}
(cont.) Show that any chief factor of a group having a composition series is a direct product of mutually isomorphic simple groups.
\end{exercise}
\begin{solution}
\end{solution}

\begin{exercise}
\label{ex:10.8}
Show that any finite group that has no proper non-trivial characteristic subgroups is a direct product of mutually isomorphic simple groups. Use this to give a new proof of Theorem~\ref{thm:10.5}.
\end{exercise}
\begin{solution}
\end{solution}

\end{document}
